\documentclass[10pt,a4paper]{amsart}
\usepackage[margin=19mm]{geometry}
\usepackage{amsmath,amssymb,mathtools}
\usepackage[scr=rsfs,cal=euler]{mathalfa}
\usepackage{xfrac}
\usepackage[unicode,hidelinks,bookmarks=false]{hyperref}
\newcommand{\ie}{i.e.\ }
\newcommand{\cf}{cf.\ }
\newcommand{\resp}{respectively\ }
\newcommand{\Subsection}{\subsection}
\providecommand{\nobreakdash}{\nobreak\hbox{-}}
\setlength{\emergencystretch}{2em}
\multlinegap0pt
\usepackage{bm}
\usepackage{bbm}
\usepackage{dsfont}
\usepackage{xspace}
\usepackage{cancel}
\usepackage{mathtools}
\usepackage{amsthm}
\makeatletter
\@ifundefined{thm}{\newtheorem{thm}{Theorem}}{}
\makeatother
\newcommand*\diff{\mathop{}\!\mathrm{d}}
\title[To report or not to report: Optimal claim reporting in a bon]{To report or not to report: Optimal claim reporting in a bonus-malus system}
\author{Lea Enzi, Stefan Thonhauser}
\date{Source: arXiv:2601.07655v1 (CC BY 4.0)}
\begin{document}
\maketitle

\noindent\emph{Source and licence.} To report or not to report: Optimal claim reporting in a bonus-malus system. Version arXiv:2601.07655v1 (CC BY 4.0), distributed under the Creative Commons Attribution 4.0 International licence, as recorded in the arXiv licence metadata for this exact version and shown on the abstract page https://arxiv.org/abs/2601.07655v1. Full rights notice: licenses/arxiv-2601.07655-CC-BY-4.0.txt.\par\smallskip

\noindent\textit{Editorial scope.} Counted unit: the Thm whose source label is \texttt{None} in the article (source0.tex lines 387--389, source label \texttt{None}), delivered together with the article's own complete original proof of it (source0.tex lines 391--566, one \texttt{proof} environment of 176 lines). No mathematics is written, repaired or re-derived: statement and proof are the article's own text line for line. Required context retained uncounted: eq:diffequation (lines 316--326). Every definition, display and label of those lines is retained unchanged. Omitted from this package and not redistributed: 1--315, 327--386, 390--390, 567--787. Nothing inside a retained excerpt is omitted or altered: every retained body is delivered exactly as the article's lines are written, so no omission receipt of any kind accompanies this package. Citations: the retained excerpts carry no citation command at all, so no bibliography entry of the article is transcribed and no reference is redistributed. Reference adaptation: none was needed and none was made; every reference inside the delivered unit resolves inside it, so no unresolved reference or citation is produced. Printed slips kept literal: no printed word, symbol, hypothesis or number of the counted statement or of its proof was altered. Layout and class: the article's own class is \texttt{scrartcl} with packages including babel, inputenc, natbib, amsmath, amsthm, amsfonts, bbm, mathtools, enumerate, xcolor, listings, algorithm2e, graphicx, subcaption; the wrapper is the lane's pinned \texttt{amsart} editorial wrapper at 10pt on A4, which restates the theorem-like environments the retained text needs and adds the article's own macro definitions that the retained text needs (\texttt{\textbackslash diff}), with extra wrapper packages (bm, bbm, dsfont, xspace, cancel, mathtools). The delivered numbering is the unit's own. Assets: no retained span contains a figure environment, a graphics-inclusion command, a PDF-page inclusion, a table or a TikZ picture; no image file is copied into this package, the package declares no asset dependency and no image file is redistributed with it. Licence: version arXiv:2601.07655v1 of this article is distributed under the Creative Commons Attribution 4.0 International licence, as recorded in the arXiv licence metadata for this exact version and shown on the abstract page https://arxiv.org/abs/2601.07655v1. A full rights notice is delivered at licenses/arxiv-2601.07655-CC-BY-4.0.txt. Characters: every retained excerpt is pure ASCII, so the delivered engine, XeLaTeX with its default Latin Modern text font, reports no missing character.
\begin{thm}
    For $i \in \{1,2\}$, $V(i,t,s,x)$ is a viscosity solution to the system of equations
    \begin{align}
        &\begin{aligned}
        &\sup_{b\geq 0} \mathcal{A}^b v_1(t,s,x)=0,\ \qquad  t \in [0,T),\ s\in [0,t],\ x \in \mathbb{R}\label{eq:diffequation}\\&\sup_{b\geq 0} \mathcal{A}^b v_2(t,s,x)=0,\ \qquad t \in [0,T),\ s\in [0,t] \text{ and } s<\mathcal{S}, \ x \in \mathbb{R} 
        \end{aligned}\\
        &v_i(T,s,x)=h(x),\\
        &v_2(t,\mathcal{S},x)=v_1(t,0,x),
    \end{align}
    where $v_i: D_i \rightarrow \mathbb{R}$.
\end{thm}

\begin{thm}
    For $j \in \{1,2\}$, let $\tilde{\underline{v}}_j:D_j\to\mathbb{R}$ and $\tilde{\overline{v}}_j:D_j\to\mathbb{R}$ be Lipschitz-continuous viscosity sub- and supersolutions to \eqref{eq:diffequation} respectively. If $\tilde{\underline{v}}_j\leq\tilde{\overline{v}}_j$ on $\partial D_j$, then $\tilde{\underline{v}}_j\leq \tilde{\overline{v}}_j$ on $D_j$. 
\end{thm}

\begin{proof}
    For $j \in \{1,2\}$, let $\tilde{\underline{v}}_j$ and $\tilde{\overline{v}}_j$ be a viscosity sub- and supersolution respectively. We argue by contradiction and assume that there is a $k \in \{1,2\}$ and a point $(t_0,x_0,s_0) \in D_k$ such that
    \begin{equation*}
        %\max_{j \in \{1,2\}} \sup_{D_j}(\underline{v}_j-\overline{v}_j)\geq
        \tilde{\underline{v}}_k(t_0,s_0,x_0)-\tilde{\overline{v}}_k(t_0,s_0,x_0)>0.
    \end{equation*}
    For $j \in \{1,2\}$, we define $\underline{v}_j(t,s,x)=e^{\delta t}\tilde{\underline{v}}_j(t,s,x)$ and $\overline{v}_j(t,s,x)=e^{\delta t}\tilde{\overline{v}}_j(t,s,x)$, with $\delta>0$. Then, these functions are viscosity sub- and supersolution respectively to the equation
    \begin{align}
    \label{proof:differentequation}
          \sup_{b\geq 0}&\left(\frac{\partial}{\partial t}+\frac{\partial}{\partial s}\right) v_j(t,s,x)+\left(c-\pi_j(s,m_j)\right)\frac{\partial}{\partial x}v_j(t,s,x) \nonumber + \lambda \int_{0}^{b}v_j(t,s,x-y)\diff F_Y(y) \nonumber \\
    &+ \lambda \int_{b}^{\infty}v_2(t,s,x-r(y,m_j))\diff F_Y(y) - (\lambda + \delta) v_j(t,s,x)=0. 
    \end{align}
    
    
    
    Let $\phi(t,x)=e^{-\gamma t}(1+x^2)$ and define $\overline{v}_j^\varepsilon(t,s,x)=\overline{v}_j(t,s,x)+\varepsilon  \phi(t,x)$. We show that there is a $\gamma>0$ such that  $\overline{v}_j^\varepsilon$ is again a supersolution to \eqref{proof:differentequation}. Therefore, for some $b\geq 0$, we consider
    \begin{align*}
       &\left(\frac{\partial}{\partial t}+\frac{\partial}{\partial s}\right) \phi(t,x)+\left(c-\pi_j(s,m_j)\right)\frac{\partial}{\partial x}\phi(t,x) \nonumber + \lambda \int_{0}^{b}\phi(t,x-y)\diff F_Y(y) \nonumber \\
    &+ \lambda \int_{b}^{\infty}\phi(t,x-r(y,m_j))\diff F_Y(y) - (\lambda+\delta) \phi(t,x)\\
    & =  e^{-\gamma t} \Big(-\gamma(1+x^2)+(c-\pi_i(s,m_j))2x + \lambda  \int_{0}^{b}(1+(x-y))^2\diff F_Y(y) \nonumber \\
    &+ \lambda \int_{b}^{\infty}(1+(x-r(y,m_j))^2)\diff F_Y(y) - (\lambda + \delta)(1+x^2)\Big)\\
    &\leq e^{-\gamma t}\Big((\lambda-\delta-\gamma)x^2 + 2(c-\pi_j(s,m_j)-\lambda\mathbb{E}[Y+r(Y,m_j)])x\\&\qquad+\lambda\mathbb{E}[Y^2+r(Y,m_j)^2]+(\lambda-\delta-\gamma) \Big).
    \end{align*}
    \\
    The expression is less than zero if the term inside the brackets is negative. Since this term is a polynomial in $x$, we can analyze its sign by studying its coefficients. If the leading coefficient is negative and the discriminant is non-positive, then the entire polynomial is non-positive for all $x$. Therefore, we need at least that $\gamma>\lambda-\delta$ and define

\begin{align*}
    &b_j(s)=2(c-\pi_j(s,m_j)-\lambda\mathbb{E}[Y+r(Y,m_j)]),\\
    &\tilde{c}_j = \lambda\mathbb{E}[Y^2+r(Y,m_j)^2].
\end{align*}
Then, the discriminant is given by
\begin{align*}
    b_j^2(s)-4(\lambda-\delta - \gamma)(\tilde{c}_j+(\lambda-\delta - \gamma)).
\end{align*}
and is less than 0 if we set 
\begin{align*}
    \gamma\geq  \max_{j \in \{1,2\}}\sup_{s \in [0,T]}\lambda-\delta + \frac{1}{2}\left(\tilde{c}_j+\sqrt{b_j(s)^2+\tilde{c}_j^2}\right).
\end{align*}
With this choice of $\gamma$ and since this holds for all $b\geq0$, we get that $\overline{v}_j^\varepsilon$ is again a supersolution to the corresponding equation. \\ \ \\
    Since $\underline{v}_j$ and $\overline{v}_j$ are Lipschitz continuous, there is a constant $L>0$ such that 
    \begin{align}
    \label{proof:growthcondition}
        \sup_{D_j}\frac{|\underline{v}_j(t,s,x)|+|\overline{v}_j(t,s,x)|}{1+|x|}<L.
    \end{align}
    Since 
    \begin{equation*}
        \underline{v}_j(t,s,x)-\overline{v}^\varepsilon_j(t,s,x)<L(1+|x|)-\varepsilon e^{-\gamma T} (1+x^2),
    \end{equation*}
    we can find a constant $a>0$ such that this is less than zero for $x>a$ and $x<-a$. Therefore, we can find a compact set $A_j \subseteq D_j$ such that
    \begin{equation*}
        M = \max_{j \in \{1,2\}} \sup_{A_j}(\underline{v}_j-\overline{v}_j^\varepsilon)= \max_{j \in \{1,2\}}\sup_{D_j}(\underline{v}_j-\overline{v}_j^\varepsilon)> \underline{v}_k(t_0,s_0,x_0)-\overline{v}_k^\varepsilon(t_0,s_0,x_0)>0,
    \end{equation*}
    for some small enough $\varepsilon>0$. We call the corresponding maximizing index $i$ and the maximizer $w^*=(t^*,s^*,x^*)$. One can show, that $\underline{v}_i$ and $\overline{v}_i^\varepsilon$ are still Lipschitz continuous on $A_i$ with constant $m$. \\
    We define,
    \begin{align*}
        g_\nu^i(t_1,s_1,x_1,t_2,s_2,x_2):=&\frac{\nu}{2}\left[(t_1-t_2)^2+(s_1-s_2)^2+(x_1-x_2)^2\right],
       % \\&+\frac{2m}{\nu^2(|t_1-t_2|+|s_1-s_2|+|x_1-x_2|+\nu},
    \end{align*}
    for some $\nu>0$ and further
    \begin{equation*}
        G_\nu ^i (t_1,s_1,x_1,t_2,s_2,x_2) := \underline{v}_i(t_1,s_1,x_1)-\overline{v}^\varepsilon_i(t_2,s_2,x_2)- g_\nu^i(t_1,s_1,x_1,t_2,s_2,x_2).
    \end{equation*}
    Let $M_\nu ^i=\max_{A_i \times A_i} G_\nu^i$ and $w_\nu=(t_1^\nu,s_1^\nu,x_1^\nu,t_2^\nu,s_2^\nu,x_2^\nu)$ the corresponding maximizer. Then,
    \begin{equation*}
        M_\nu ^i \geq G_\nu^i(t^*,s^*,x^*,t^*,s^*,x^*)=M.
    \end{equation*}
    We need that $(t_1^\nu,s_1^\nu,x_1^\nu,t_2^\nu,s_2^\nu,x_2^\nu) \notin \partial (A_i \times A_i)$. %, which we will discuss that later. 
    The boundaries are of the following form: 
   \begin{align*}
\partial (A_1 \times A_1) =  \big\{\, &(t_1, t_2, s_1, s_2, x_1, x_2) \in A_1 \ \big| \\
& t_j \in \{0,T\} \ \text{or} \ s_j \in \{0,T\} \ \text{or} \ s_j = t_j \ \text{or} \ x_j \in \{-a,a\} \text{ for } j \in \{1,2\}\big\}.
\end{align*}
\begin{align*}
\partial (A_2 \times A_2) =  \big\{\, &(t_1, t_2, s_1, s_2, x_1, x_2) \in A_2 \ \big| \\
& t_j \in \{0,T\} \ \text{or} \ s_j \in \{0,\mathcal{S}\} \ \text{or} \ s_j = t_j \ \text{or} \ x_j \in \{-a,a\} \text{ for } j \in \{1,2\}\big\}.
\end{align*}
    By \eqref{proof:growthcondition} and the construction of $A_i$, we know that
    \begin{equation*}
         G_\nu ^i (t_1,s_1,x_1,t_2,s_2,x_2) \leq 2L (1+a)- \frac{\nu}{2}\left[(t_1-t_2)^2+(s_1-s_2)^2+(x_1-x_2)^2\right].
    \end{equation*}
    If either $|t_1-t_2|>\vartheta$, $|s_1-s_2|>\vartheta$ or $|x_1-x_2|>\vartheta$ for some $\vartheta>0$, then this is negative for $\nu > \frac{4L(1+a)}{\vartheta}$. \\So let us consider the case where $t_1=t_2$, $s_1=s_2$ and $x_1=x_2$.
    Then,
    \begin{align*}
             G_\nu ^i (t_1,s_1,x_1,t_1,s_1,x_1) =  \underline{v}_i(t_1,s_1,x_1)-\overline{v}^\varepsilon_i(t_1,s_1,x_1).
    \end{align*}
    In the cases where $t_1=T$,  $s_1=\mathcal{S}$ (if $i=2$), $t_1=0$, $s_1=0$ or $t_1=s_1$, this is less than or equal to zero by the theorem's assumption. \\If $x_1 \in \{-a,a\}$, this is less than or equal to zero by construction of $A_i$. Therefore, the maximum is not attained on the boundary. \\ \ \\
    \noindent
    Since $G_\nu^i$ attains a maximum in $w_\nu$, it holds that for any $(t_1,s_1,x_1) \in D_i$, 
    \begin{equation*}
        G_\nu^i(t_1,s_1,x_1,t_2^\nu,s_2^\nu,x_2^\nu)\leq G_\nu^i(w_\nu),
    \end{equation*}
    and therefore
    \begin{equation*}
        \underline{v}_i(t_1,s_1,x_1)-\underline{v}_i(t_1^\nu,s_1^\nu,x_1^\nu)\leq g_\nu^i(w_\nu)-g_\nu^i (t_1,s_1,x_1,t_2^\nu,s_2^\nu,x_2^\nu).
    \end{equation*}
    Together with Taylor's formula we arrive at
    \begin{align*}&\limsup_{(t_1,s_1,x_1)\rightarrow (t_1^\nu,s_1^\nu,x_1^\nu)}\biggl(\frac{ \underline{v}_i(t_1,s_1,x_1)-\underline{v}_i(t_1^\nu,s_1^\nu,x_1^\nu)}{|t_1^\nu-t_1|+|s_1^\nu-s_1|+|x_1^\nu-x_1|}\\
        &-\frac{ (t_1^\nu-t_1)\frac{\partial}{\partial t_1}g_\nu(w_\nu)+(s_1^\nu-s_1)\frac{\partial}{\partial s_1}g_\nu(w_\nu)+(x_1^\nu-x_1)\frac{\partial}{\partial x_1}g_\nu(w_\nu)}{|t_1^\nu-t_1|+|s_1^\nu-s_1|+|x_1^\nu-x_1|}\biggr)\leq 0
    \end{align*}
    and therefore $\partial_{t_1+s_1+x_1} g_\nu^i(w_\nu) \in D^+ (\underline{v}_i)(t_1^\nu,s_1^\nu,x_1^\nu)$, which is the set of all superdifferentials of $\underline{v}_i$ at $(t_1^\nu,s_1^\nu,x_1^\nu)$.
    With a similar argument,
     \begin{align*}
        &\liminf_{(t_2,s_2,x_2)\rightarrow (t_2^\nu,s_2^\nu,x_2^\nu)} \biggl(\frac{ \overline{v}^\varepsilon_i(t_2,s_2,x_2)-\overline{v}^\varepsilon_i(t_2^\nu,s_2^\nu,x_2^\nu)}{|t_2^\nu-t_2|+|s_2^\nu-s_2|+|x_2^\nu-x_2|}\\
        &-\frac{(t_2^\nu-t_2)(-\frac{\partial}{\partial t_2}g_\nu(w_\nu))+(s_2^\nu-s_2)(-\frac{\partial}{\partial s_2}g_\nu(w_\nu))+(x_2^\nu-x_2)(-\frac{\partial}{\partial x_2}g_\nu(w_\nu))}{|t_2^\nu-t_2|+|s_2^\nu-s_2|+|x_2^\nu-x_2|}\biggr)\geq 0.
    \end{align*}
    Therefore, $-\partial_{t_2+s_2+x_2} g_\nu^i(w_\nu) \in D^- (\overline{v}^\varepsilon_i)(t_2^\nu,s_2^\nu,x_2^\nu)$, which is the set of all subdifferentials of $\overline{v}^\varepsilon_i$ in $(t_2^\nu,s_2^\nu,x_2^\nu)$. In addition, 
    \begin{equation}
    \label{eq:g_x=-g_y}
        \partial_{t_1+s_1+x_1} g_\nu^i(w_\nu)=-\partial_{t_2+s_2+x_2} g_\nu^i(w_\nu).
    \end{equation}
    
    \noindent
    Now we return to the equation \eqref{proof:differentequation}. There is a $b\geq 0$ such that
    \begin{align}
    \label{eq:proof>=0}
        &\left(\frac{\partial}{\partial t_1}+\frac{\partial}{\partial s_1}\right) g_{\nu}^i(t_1,s_1,x_1,t_2,s_2,x_2)|_{w_{\nu}} \nonumber \\&+\left(c-\pi_i(s_1^\nu,m_i)\right)\frac{\partial}{\partial x_1}g_{\nu}^i(t_1,s_1,x_1,t_2,s_2,x_2)|_{w_{\nu}} \nonumber \\
    &+ \lambda \left(\int_{0}^{b}\underline{v}_i(t_1^\nu,s_1^\nu,x_1^\nu-y)\diff F_Y(y) \nonumber 
    +\int_{b}^{\infty}\underline{v}_2(t_1^\nu,0,x_1^\nu-r(y,m_i))\diff F_Y(y)\right) \\&- (\lambda+\delta) \underline{v}_i(t_1^\nu,s_1^\nu,x_1^\nu) \geq 0,
    \end{align}
    and 
    \begin{align}
    \label{eq:proof<=0}
        &-\left(\frac{\partial}{\partial t_2}+\frac{\partial}{\partial s_2}\right) g_{\nu}^i(t_1,s_1,x_1,t_2,s_2,x_2)|_{w_{\nu}}\\
        &-\left(c-\pi_i(s_2^\nu,m_i)\right)\frac{\partial}{\partial x_2}g_{\nu}^i(t_1,s_1,x_1,t_2,s_2,x_2)|_{w_{\nu}} \nonumber \\
    &+ \lambda \left(\int_{0}^{b}\overline{v}^\varepsilon_i(t_2^\nu,s_2^\nu,x_2^\nu-y)\diff F_Y(y) \nonumber 
    +\int_{b}^{\infty}\overline{v}^\varepsilon_2(t_2^\nu,0,x_2^\nu-r(y,m_i))\diff F_Y(y)\right) \\&- (\lambda+\delta) \overline{v}^\varepsilon_i(t_2^\nu,s_2^\nu,x_2^\nu) \leq 0,
    \end{align}
    By subtracting the left hand side of \eqref{eq:proof>=0} from the LHS of \eqref{eq:proof<=0} together with \eqref{eq:g_x=-g_y}, we arrive at
  
    \begin{align}
    \label{proof:(delta+lambda)}
         &(\lambda + \delta)(\underline{v}_i(t_1^\nu,s_1^\nu,x_1^\nu)-\overline{v}^\varepsilon_i(t_2^\nu,s_2^\nu,x_2^\nu)) \nonumber\\
        &\leq \left(c-\pi_i(s_2^\nu,m_i)\right)\frac{\partial}{\partial x_2}g_{\nu}^i(t_1,s_1,x_1,t_2,s_2,x_2)|_{w_{\nu}} \nonumber \\&+ \left(c-\pi_i(s_1^\nu,m_i)\right)\frac{\partial}{\partial x_1}g_{\nu}^i(t_1,s_1,x_1,t_2,s_2,x_2)|_{w_{\nu}}\nonumber\\
        &+\lambda \left(\int_{0}^{b}\underline{v}_i(t_1^\nu,s_1^\nu,x_1^\nu-y)\diff F_Y(y)  
   +\int_{b}^{\infty}\underline{v}_2(t_1^\nu,0,x_1^\nu-r(y,m_i))\diff F_Y(y)\right) \nonumber\\
   &-  \lambda \left(\int_{0}^{b}\overline{v}^\varepsilon_i(t_2^\nu,s_2^\nu,x_2^\nu-y)\diff F_Y(y)  
    +\int_{b}^{\infty}\overline{v}^\varepsilon_2(t_2^\nu,0,x_2^\nu-r(y,m_i))\diff F_Y(y)\right)
    \end{align}
    Since $w_\nu$ is the maximum of $G_\nu$ on $A_i \times A_i$,  we have that 
    \begin{align*}
        G_\nu ^i (t_1^\nu,s_1^\nu,x_1^\nu,t_1^\nu,s_1^\nu,x_1^\nu)+G_\nu ^i (t_2^\nu,s_2^\nu,x_2^\nu,t_2^\nu,s_2^\nu,x_2^\nu) \leq 2G_\nu(w_\nu).
    \end{align*}
    This is equivalent to
    \begin{align*}
        &\underline{v}_i(t_1^\nu,s_1^\nu,x_1^\nu)-\overline{v}_i^\varepsilon(t_1^\nu,s_1^\nu,x_1^\nu)+ \underline{v}_i(t_2^\nu,s_2^\nu,x_2^\nu)-\overline{v}_i^\varepsilon(t_2^\nu,s_2^\nu,x_2^\nu) \\
        &\leq 2(\underline{v}_i(t_1^\nu,s_1^\nu,x_1^\nu)-\overline{v}_i^\varepsilon(t_2^\nu,s_2^\nu,x_2^\nu)) - \nu ||(t_1^\nu,s_1^\nu,x_1^\nu)-(t_2^\nu,s_2^\nu,x_2^\nu)||^2_2
    \end{align*}
    and further 
    \begin{align*}
       &\nu||(t_1^\nu,s_1^\nu,x_1^\nu)-(t_2^\nu,s_2^\nu,x_2^\nu)||^2_2 \\
       &\leq \underline{v}_i(t_1^\nu,s_1^\nu,x_1^\nu)- \underline{v}_i(t_2^\nu,s_2^\nu,x_2^\nu)+ \overline{v}_i^\varepsilon(t_1^\nu,s_1^\nu,x_1^\nu)-\overline{v}_i^\varepsilon(t_2^\nu,s_2^\nu,x_2^\nu)\\
       & \leq m ||(t_1^\nu,s_1^\nu,x_1^\nu)-(t_2^\nu,s_2^\nu,x_2^\nu)||_1
       \leq m \sqrt{2} ||(t_1^\nu,s_1^\nu,x_1^\nu)-(t_2^\nu,s_2^\nu,x_2^\nu)||_2.
    \end{align*}
    Thus,
    \begin{equation}
    \label{proof:inequality}
        ||(t_1^\nu,s_1^\nu,x_1^\nu)-(t_2^\nu,s_2^\nu,x_2^\nu)||_2 \leq \frac{\sqrt{2}m}{\nu}.
    \end{equation}
    So if we take a sequence $(\nu_n)_{n \in \mathbb{N}}$ such that $w_\nu$ converges to $(\overline{t}_1,\overline{s}_1,\overline{x}_1,\overline{t}_2,\overline{s}_2,\overline{x}_2)$ as $\nu_n \rightarrow \infty$, then we get by \eqref{proof:inequality} that $(\overline{t}_1,\overline{s}_1,\overline{x}_1)=(\overline{t}_2,\overline{s}_2,\overline{x}_2)$.

By \eqref{proof:(delta+lambda)} and \eqref{eq:g_x=-g_y}, we arrive at
   \begin{align*}
         &(\lambda + \delta)(\underline{v}_i(\overline{t}_1,\overline{s}_1,\overline{x}_1)-\overline{v}^\varepsilon_i(\overline{t}_1,\overline{s}_1,\overline{x}_1)) \nonumber\\
        &\leq \lambda \left(\int_{0}^{b}\underline{v}_i(\overline{t}_1,\overline{s}_1,\overline{x}_1-y)\diff F_Y(y)  
   +\int_{b}^{\infty}\underline{v}_2(\overline{t}_1,0,\overline{x}_1-r(y,m_i))\diff F_Y(y)\right) \nonumber\\
   &-  \lambda \left(\int_{0}^{b}\overline{v}^\varepsilon_i(\overline{t}_1,\overline{s}_1,\overline{x}_1-y)\diff F_Y(y)  
    +\int_{b}^{\infty}\overline{v}^\varepsilon_2(\overline{t}_1,0,\overline{x}_1-r(y,m_i))\diff F_Y(y)\right) \leq \lambda M.
    \end{align*}
    So 
    \begin{equation*}
       M \leq \lim_{n \to \infty}M_{\nu_n} =\underline{v}_i(\overline{t}_1,\overline{s}_1,\overline{x}_1)-\overline{v}^\varepsilon_i(\overline{t}_1,\overline{s}_1,\overline{x}_1)\leq \frac{\lambda}{\lambda + \delta} M,
    \end{equation*}
    which is a contradiction and thus completes the proof.

\end{proof}

\end{document}
